\noindent\textbf{ON NETWORK ADJUSTMENT.}\hfill323
\begin{center}\textsc{Gauss} to \textsc{Gerling}.\quad Göttingen, 2 May 1837.\end{center}

.\ .\ .\ .\ .\ .\ As for the computation of the geodetic measurements, I did carry
out the elimination for my own principal triangles from Göttingen to Jever,
about 40 triangles between 33 points, completely and with full rigor. Even
with the devices I did in fact use, it was a very arduous and protracted
task at the time. Merely to explain these devices to you adequately would
take a small book, and to write it I would first have to study the matter
afresh, which I cannot contemplate now.

\begingroup
\def\letterfourstars{*}
\makeatletter
\renewcommand{\@makefnmark}{\hbox{\normalfont[\textsuperscript{\letterfourstars})]}}
\renewcommand{\@makefntext}[1]{\parindent1em\noindent\hbox to2em{\hss[\textsuperscript{\letterfourstars})\ }#1]}
\makeatother
For the later triangles surveyed by my son, \textsc{Hartmann}, and \textsc{Müller}
in Westphalia, the Lüneburg region, the Harz, and the Weser region, as well
as for last year's triangles, I relaxed the utmost rigor somewhat and
proceeded as follows.
\begin{list}{}{\setlength{\leftmargin}{2.4em}\setlength{\labelwidth}{1.5em}\setlength{\labelsep}{.5em}\setlength{\itemsep}{.3em}\setlength{\topsep}{.3em}}
\item[1)] I considered the angle conditions \textls[120]{alone} and adjusted rigorously on that basis. Then
\item[2)] I sought the new adjustment to these adjusted numbers required by the
ratios of the sides\footnote{See p.\ 253.}.
\end{list}
\def\letterfourstars{**}
\par\noindent\hspace*{\parindent}If one takes the trouble to apply this procedure alternately and repeatedly
until the results stabilize, as I already recommended in
Supplementum etc.\ p.\ 28\footnote{Volume IV, pp.\ 77/79.}, one arrives, as proved there, at
\textls[120]{exactly} the same results as if one took all the conditions into account at once.
\par\medskip
.\ .\ .\ .\ .\ .
\endgroup
